% GENERATED FILE. Edit canonical lessons, then run npm run generate:books.
% canonical-source-hash: fnv1a64:81e66395016a3993
% canonical-lessons: JA-C132-soko, JA-W132-so, JA-C132-soto, JA-C132-kore, JA-W132-re, JA-C132-sore, JA-C132-kuruma, JA-W132-ru, JA-R132-this-and-that

\chapter{This and That: So, Re and Ru}
\label{ch:ja-this-and-that-so-re-ru}
\hlchaptermodality{132}

\begin{chapteropening}
\textbf{By the end of this chapter:} \emph{I can write so, re and ru, read words that use them, and point to things and places near me and near the listener.}

Read soko, soto, kore, sore and kuruma, point with kore and sore, write kuruma from memory with \ja{る} in the middle, and write the three new signs once each.
\end{chapteropening}

\section[\texorpdfstring{\ja{そこ} (soko)}{そこ (soko)}]{\ja{そこ} (soko) — there, near you}

\label{lesson:JA-C132-soko}



\begin{quote}
Before the new word, say \emph{koko}, here. \emph{Write it:} \textbf{\ja{こ}}; then \textbf{\ja{ょ}} — \textbf{R2}, five lessons back

\begin{itemize}
\raggedright
  \item \emph{Recall:} say \emph{a Japanese serow} — \textbf{R4}, eighty lessons back
  \item \emph{Recall:} say \emph{persistent} — \textbf{R3}, twenty lessons back
\end{itemize}
\end{quote}

\subsection*{You'll want to know: \ja{そこ}}
\textbf{\ja{そこ}} — \emph{soko} — \textquotedblleft{}there\textquotedblright{}, the place near the person you are talking to.

You know \textbf{\ja{ここ}}, \emph{koko}, here, the place near you. Change the first sign and the place moves across to the listener: \emph{soko}, there, by them. The second sign is the \textbf{\ja{こ}} you already write. The first one is new, and the next lesson writes it.

\subsection*{Your turn}
\begin{itemize}
\raggedright
  \item \emph{Say it:} \emph{soko}
  \item \emph{Say it:} \emph{koko}, picturing the spot at your own feet, then \emph{soko}, picturing the spot at the listener's
  \item \emph{Recall:} name the sign in \textbf{\ja{そこ}} you can already write, and the one you cannot
\end{itemize}

\begin{tcolorbox}[breakable,colback=teal!4,colframe=teal!35!black,title={Before you move on}]
What is \textquotedblleft{}there\textquotedblright{}, near the listener? (\textbf{\ja{そこ}}, \emph{soko}.) And \textquotedblleft{}here\textquotedblright{}? (\textbf{\ja{ここ}}, \emph{koko}.)
\end{tcolorbox}
\section[\texorpdfstring{\ja{そ} (so)}{そ (so)}]{\ja{そ} — one stroke, two zigzags}

\label{lesson:JA-W132-so}



\begin{quote}
Four recalls, then the sign.

\begin{itemize}
\raggedright
  \item \emph{Recall:} say \emph{there}, near the listener — \textbf{R1}, one lesson back
  \item \emph{Recall:} say \emph{a library} — \textbf{R2}, five lessons back
  \item \emph{Recall:} say \emph{a water flask} — \textbf{R4}, eighty lessons back
  \item \emph{Recall:} say \emph{noisy} — \textbf{R3}, twenty lessons back
\end{itemize}
\end{quote}

\subsection*{Script — the shape}
\hlblockfigure{\detokenize{figures/JA-W132-so-filmstrip.pdf}}{How \ja{そ} is written, stroke by stroke}

\begin{quote}
\textbf{\ja{そ}} — \emph{so}
\end{quote}

One stroke, and the pen never lifts:

\begin{enumerate}
\raggedright
  \item a short \textbf{bar} at the top, left to right
  \item turn sharply and run \textbf{down and to the left} until you reach the long middle bar
  \item turn back and run \textbf{right} along the whole of that bar
  \item double back a little way, then drop into an \textbf{open curve} that rounds the lower left and finishes along the base to the right
\end{enumerate}

The pen zigzags twice before it reaches the curve. Some typefaces print the top apart from the rest; written by hand here, it is one unbroken stroke.

With \textbf{\ja{こ}} after it, you can now write the word from the last lesson: \textbf{\ja{そこ}}, \emph{soko}, there.

\subsection*{Your turn}
\begin{enumerate}
\raggedright
  \item Trace \textbf{\ja{そ}} saying \textquotedblleft{}bar, down, across, curve\textquotedblright{}.
  \item Write \textbf{\ja{そこ}} from memory and say \emph{soko}.
  \item Write \textbf{\ja{ここ}} and \textbf{\ja{そこ}} side by side, and say which one is the place by the listener.
\end{enumerate}

\begin{tcolorbox}[breakable,colback=teal!4,colframe=teal!35!black,title={Before you move on}]
Stop after one accurate form.

Source: \href{https://www.unicode.org/charts/PDF/U3040.pdf}{Unicode Hiragana chart}; stroke count and order per \href{https://kanjivg.tagaini.net/kanjivg/kanji/0305d.html}{KanjiVG}.
\end{tcolorbox}
\section[\texorpdfstring{\ja{そと} (soto)}{そと (soto)}]{\ja{そと} (soto) — outside}

\label{lesson:JA-C132-soto}



\begin{quote}
\emph{Write it:} \textbf{\ja{そ}} — \textbf{R1}, one lesson back

Then ask for tea politely — \textbf{R2}, five lessons back.

\begin{itemize}
\raggedright
  \item \emph{Recall:} say \emph{a rope of straw} — \textbf{R4}, eighty lessons back
  \item \emph{Recall:} say \emph{shameless} — \textbf{R3}, twenty lessons back
\end{itemize}
\end{quote}

\subsection*{You'll want to know: \ja{そと}}
\textbf{\ja{そと}} — \emph{soto} — \textquotedblleft{}outside\textquotedblright{}.

The new sign is spent at once. \textbf{\ja{そ}} with \textbf{\ja{と}} after it: two beats, \emph{so–to}, and both signs are ones you write. \textbf{\ja{そこ}} is a spot by the listener; \textbf{\ja{そと}} is anywhere beyond the walls, out of doors.

\subsection*{Your turn}
\begin{itemize}
\raggedright
  \item \emph{Say it:} \emph{soto}
  \item \emph{Say it:} \emph{soto}, picturing a window or a door
  \item \emph{Write it:} \textbf{\ja{そと}} from memory, one stroke for the first sign and two for the second
\end{itemize}

\begin{tcolorbox}[breakable,colback=teal!4,colframe=teal!35!black,title={Before you move on}]
What does \ja{そと} mean? (\textquotedblleft{}Outside\textquotedblright{}.) Say it once more.
\end{tcolorbox}
\section[\texorpdfstring{\ja{これ} (kore)}{これ (kore)}]{\ja{これ} (kore) — this}

\label{lesson:JA-C132-kore}



\begin{quote}
Say \emph{outside} — \textbf{R1}, one lesson back. \emph{Write it:} \textbf{\ja{を}} — \textbf{R2}, five lessons back

\begin{itemize}
\raggedright
  \item \emph{Recall:} say \emph{a ladder} — \textbf{R4}, eighty lessons back
  \item \emph{Recall:} say \emph{reliable} — \textbf{R3}, twenty lessons back
\end{itemize}
\end{quote}

\subsection*{You'll want to know: \ja{これ}}
\textbf{\ja{これ}} — \emph{kore} — \textquotedblleft{}this\textquotedblright{}, a thing near you as you speak.

\textbf{\ja{ここ}} is a place near you; \textbf{\ja{これ}} is a thing near you, the one in your hand or under your finger. It begins with the same \textbf{\ja{こ}}. The second sign is one you do not write yet. It turns the pointing word into a thing, and the next lesson writes it.

\subsection*{Your turn}
\begin{itemize}
\raggedright
  \item \emph{Say it:} \emph{kore}
  \item \emph{Say it:} \emph{kore}, picturing a pen in your hand, then \emph{koko}, picturing the spot where you stand
  \item \emph{Recall:} name the sign in \textbf{\ja{これ}} you can already write, and the one you cannot
\end{itemize}

\begin{tcolorbox}[breakable,colback=teal!4,colframe=teal!35!black,title={Before you move on}]
What is \textquotedblleft{}this\textquotedblright{}? (\textbf{\ja{これ}}, \emph{kore}.) And \textquotedblleft{}here\textquotedblright{}? (\textbf{\ja{ここ}}, \emph{koko}.)
\end{tcolorbox}
\section[\texorpdfstring{\ja{れ} (re)}{れ (re)}]{\ja{れ} — one first stroke, a third ending}

\label{lesson:JA-W132-re}



\begin{quote}
Four recalls, then the sign.

\begin{itemize}
\raggedright
  \item \emph{Recall:} say \emph{this}, a thing near you — \textbf{R1}, one lesson back
  \item \emph{Recall:} write \textbf{\ja{わ}} and \textbf{\ja{ね}}, and say each one
  \item \emph{Recall:} say \emph{a chain} — \textbf{R4}, eighty lessons back
  \item \emph{Recall:} say \emph{pure white} — \textbf{R3}, twenty lessons back
\end{itemize}
\end{quote}

\subsection*{Script — the shape}
\hlblockfigure{\detokenize{figures/JA-W132-re-filmstrip.pdf}}{How \ja{れ} is written, stroke by stroke}

\begin{quote}
\textbf{\ja{れ}} — \emph{re}
\end{quote}

Two strokes, one lift:

\begin{enumerate}
\raggedright
  \item the long \textbf{vertical}, from top to bottom
  \item a short \textbf{bar} from the left across the vertical; without lifting, turn sharply and cut \textbf{down and to the left} through the vertical, climb back up that line and rise to the right over an \textbf{arch}, come down its right side, round the foot and \textbf{flick} up to the right
\end{enumerate}

The first stroke and the zigzag are the ones your hand already makes in \textbf{\ja{わ}} and \textbf{\ja{ね}}. Only the ending of the second stroke is new:

\noindent
\begin{tabularx}{\linewidth}{@{}>{\raggedright\arraybackslash}X>{\raggedright\arraybackslash}X@{}}
\toprule
\textbf{sign} & \textbf{how the second stroke ends} \\
\midrule
\textbf{\ja{わ}} & a broad belly that curves back to the left \\
\textbf{\ja{ね}} & a small loop at the lower right \\
\textbf{\ja{れ}} & a drop and a flick out to the right \\
\bottomrule
\end{tabularx}

With \textbf{\ja{こ}} in front of it, you can now write the word from the last lesson: \textbf{\ja{これ}}, \emph{kore}, this.

\subsection*{Your turn}
\begin{enumerate}
\raggedright
  \item Trace \textbf{\ja{れ}} saying \textquotedblleft{}down; bar, zigzag, arch, flick\textquotedblright{}.
  \item Write \textbf{\ja{わ}}, \textbf{\ja{ね}} and \textbf{\ja{れ}} in a row, and look only at the three endings.
  \item Write \textbf{\ja{これ}} from memory and say \emph{kore}.
\end{enumerate}

\begin{tcolorbox}[breakable,colback=teal!4,colframe=teal!35!black,title={Before you move on}]
Stop after one accurate form.

Source: \href{https://www.unicode.org/charts/PDF/U3040.pdf}{Unicode Hiragana chart}; stroke count and order per \href{https://kanjivg.tagaini.net/kanjivg/kanji/0308c.html}{KanjiVG}.
\end{tcolorbox}
\section[\texorpdfstring{\ja{それ} (sore)}{それ (sore)}]{\ja{それ} (sore) — that, near you}

\label{lesson:JA-C132-sore}



\begin{quote}
\emph{Write it:} \textbf{\ja{れ}} — \textbf{R1}, one lesson back

Then say \emph{there}, near the listener — \textbf{R2}, five lessons back.

\begin{itemize}
\raggedright
  \item \emph{Recall:} say \emph{cotton wool} — \textbf{R4}, eighty lessons back
  \item \emph{Recall:} say \emph{pitch black} — \textbf{R3}, twenty lessons back
\end{itemize}
\end{quote}

\subsection*{You'll want to know: \ja{それ}}
\textbf{\ja{それ}} — \emph{sore} — \textquotedblleft{}that\textquotedblright{}, a thing near the person you are talking to.

You have met both halves. \textbf{\ja{そ}} points at the listener, as in \textbf{\ja{そこ}}; \textbf{\ja{れ}} makes it a thing, as in \textbf{\ja{これ}}. Put side by side, the four words line up:

\noindent
\begin{tabularx}{\linewidth}{@{}>{\raggedright\arraybackslash}X>{\raggedright\arraybackslash}X@{}}
\toprule
\textbf{near you} & \textbf{near the listener} \\
\midrule
\textbf{\ja{ここ}}, here & \textbf{\ja{そこ}}, there \\
\textbf{\ja{これ}}, this & \textbf{\ja{それ}}, that \\
\bottomrule
\end{tabularx}

\subsection*{Your turn}
\begin{itemize}
\raggedright
  \item \emph{Say it:} \emph{sore}
  \item \emph{Say it:} \emph{kore}, picturing something in front of you, then \emph{sore}, picturing something by the listener
  \item \emph{Write it:} \textbf{\ja{それ}} from memory; you write every sign in it now
\end{itemize}

\begin{tcolorbox}[breakable,colback=teal!4,colframe=teal!35!black,title={Before you move on}]
What does \ja{それ} mean? (\textquotedblleft{}That\textquotedblright{}, near the listener.) And \ja{これ}? (\textquotedblleft{}This\textquotedblright{}.)
\end{tcolorbox}
\section[\texorpdfstring{\ja{くるま} (kuruma)}{くるま (kuruma)}]{\ja{くるま} (kuruma) — a car}

\label{lesson:JA-C132-kuruma}



\begin{quote}
Say \emph{that}, a thing near the listener — \textbf{R1}, one lesson back. \emph{Write it:} \textbf{\ja{そ}} — \textbf{R2}, five lessons back

\begin{itemize}
\raggedright
  \item \emph{Recall:} say \emph{a bookshop} — \textbf{R4}, eighty lessons back
  \item \emph{Recall:} say \emph{bright red} — \textbf{R3}, twenty lessons back
\end{itemize}
\end{quote}

\subsection*{You'll want to know: \ja{くるま}}
\textbf{\ja{くるま}} — \emph{kuruma} — \textquotedblleft{}a car\textquotedblright{}.

Three beats: \emph{ku–ru–ma}. The first and last signs are \textbf{\ja{く}} and \textbf{\ja{ま}}, and you write both. The middle one is new. It looks like \textbf{\ja{ろ}} with a small loop tied at the bottom, and the next lesson writes it.

\subsection*{Your turn}
\begin{itemize}
\raggedright
  \item \emph{Say it:} \emph{kuruma}
  \item \emph{Say it:} \emph{kuruma}, then count its beats aloud — three
  \item \emph{Recall:} name the sign in \textbf{\ja{くるま}} you cannot write yet, and say which sign it looks like
\end{itemize}

\begin{tcolorbox}[breakable,colback=teal!4,colframe=teal!35!black,title={Before you move on}]
What is \textquotedblleft{}a car\textquotedblright{}? (\textbf{\ja{くるま}}, \emph{kuruma}.) How many beats? (\textbf{Three}.)
\end{tcolorbox}
\section[\texorpdfstring{\ja{る} (ru)}{る (ru)}]{\ja{る} — \ja{ろ}, with a loop}

\label{lesson:JA-W132-ru}



\begin{quote}
Five recalls, then the sign.

\begin{itemize}
\raggedright
  \item \emph{Recall:} say \emph{a car}, and clap its beats — \textbf{R1}, one lesson back
  \item \emph{Recall:} say \emph{outside} — \textbf{R2}, five lessons back
  \item \emph{Recall:} write \textbf{\ja{ろ}}: the bar, the diagonal, the open belly
  \item \emph{Recall:} say \emph{a florist} — \textbf{R4}, eighty lessons back
  \item \emph{Recall:} say \emph{vague} — \textbf{R3}, twenty lessons back
\end{itemize}
\end{quote}

\subsection*{Script — the shape}
\hlblockfigure{\detokenize{figures/JA-W132-ru-filmstrip.pdf}}{How \ja{る} is written, stroke by stroke}

\begin{quote}
\textbf{\ja{る}} — \emph{ru}
\end{quote}

One stroke, and the pen never lifts:

\begin{enumerate}
\raggedright
  \item a short \textbf{bar} at the top, left to right
  \item turn sharply and run the long \textbf{diagonal} down to the lower left
  \item swing back up to the right and round the big \textbf{bowl}
  \item at the bottom, curl up into a small \textbf{loop} and finish where it meets the base
\end{enumerate}

Up to the bowl, this is \textbf{\ja{ろ}}. Only the ending differs:

\noindent
\begin{tabularx}{\linewidth}{@{}>{\raggedright\arraybackslash}X>{\raggedright\arraybackslash}X@{}}
\toprule
\textbf{you write} & \textbf{it ends with} \\
\midrule
\textbf{\ja{ろ}}, \emph{ro} & a short tail curving left \\
\textbf{\ja{る}}, \emph{ru} & a small closed loop \\
\bottomrule
\end{tabularx}

With \textbf{\ja{く}} and \textbf{\ja{ま}} around it, you can now write the word from the last lesson: \textbf{\ja{くるま}}, \emph{kuruma}, a car.

\subsection*{Your turn}
\begin{enumerate}
\raggedright
  \item Trace \textbf{\ja{る}} saying \textquotedblleft{}bar, diagonal, bowl, loop\textquotedblright{}.
  \item Write \textbf{\ja{ろ}} and \textbf{\ja{る}} side by side; check that only the second one closes a loop.
  \item Write \textbf{\ja{くるま}} from memory and say \emph{kuruma}.
\end{enumerate}

\begin{tcolorbox}[breakable,colback=teal!4,colframe=teal!35!black,title={Before you move on}]
Stop after one accurate form.

Source: \href{https://www.unicode.org/charts/PDF/U3040.pdf}{Unicode Hiragana chart}; stroke count and order per \href{https://kanjivg.tagaini.net/kanjivg/kanji/0308b.html}{KanjiVG}.
\end{tcolorbox}
\section[(dialogue)]{This and that — the three signs, written}

\label{lesson:JA-R132-this-and-that}



\begin{quote}
\emph{Write it:} \textbf{\ja{る}} — \textbf{R1}, one lesson back

Then say \emph{this} — \textbf{R2}, five lessons back.

\begin{itemize}
\raggedright
  \item \emph{Recall:} say \emph{a fishmonger} — \textbf{R4}, eighty lessons back
  \item \emph{Recall:} say \emph{certain} — \textbf{R3}, twenty lessons back
\end{itemize}
\end{quote}

\subsection*{Your turn — read the five words}
\emph{Read it:} each one below aloud, then say what it means

Read these aloud:

\begin{itemize}
\raggedright
  \item \textbf{\ja{そこ}} — there, near the listener
  \item \textbf{\ja{そと}} — outside
  \item \textbf{\ja{これ}} — this, near you
  \item \textbf{\ja{それ}} — that, near the listener
  \item \textbf{\ja{くるま}} — a car
\end{itemize}

Two signs do the pointing. \textbf{\ja{こ}} points near you and \textbf{\ja{そ}} points near the listener; the sign after them says whether you mean a place or a thing.

\subsection*{Your turn — point and write}
\begin{enumerate}
\raggedright
  \item \emph{Write it:} \textbf{\ja{これ}} and \textbf{\ja{それ}} from memory, and circle the sign that changes
  \item \emph{Write it:} \textbf{\ja{そ}}, \textbf{\ja{れ}} and \textbf{\ja{る}} once each
  \item Picture something in front of you and say \emph{kore}; picture something by the listener and say \emph{sore}; picture a window and say \emph{soto}. \emph{Write it:} \textbf{\ja{くるま}}
\end{enumerate}

\begin{tcolorbox}[breakable,colback=teal!4,colframe=teal!35!black,title={Before you move on}]
There? (\textbf{\ja{そこ}}, \emph{soko}.) Outside? (\textbf{\ja{そと}}, \emph{soto}.) This? (\textbf{\ja{これ}}, \emph{kore}.) That? (\textbf{\ja{それ}}, \emph{sore}.) A car? (\textbf{\ja{くるま}}, \emph{kuruma}.)
\end{tcolorbox}
